\documentclass[11pt]{article}
\usepackage[margin=1in]{geometry}
\usepackage{amsmath}
\usepackage{amssymb}

\setlength{\parindent}{0pt}

\begin{document}

\textbf{Question 5.2}
\bigskip

What is the present value of a security that will pay \$29,000 in 20 years if securities of equal risk pay 5\% annually?

\bigskip
\textbf{Answer}

\[
  29000 = PV(1.05)^{20}
\]
\[
  \Rightarrow \boxed {PV \approx 10929.795}
\]

\bigskip
\bigskip
\textbf{Question 5.5}
\bigskip

You have \$33,556.25 in a brokerage account, and to deposit an additional \$50,000 at the end of every future year until your account totals \$220,000.
You expect to earn 12\% annually on the account. How many years will it take to reach your goal?

\bigskip
\textbf{Answer}

\[
  220000 = 33556.25(1.12)^n + 50000[\frac{1.12^n - 1}{0.12}]
\]

Let $x=1.12^n$:

\[
  33556.25x+416666.67(x - 1)=220000
\]
\[
  \Rightarrow 450222.92x=636666.67
\]
\[
  x=1.4142
\]
\[
  \because 1.4142 = 1.12^n \therefore n = \frac{\ln(1.4142)}{\ln(1.12)}
\]
\[
  \boxed{n \approx 3.06 \ years}
\]

\bigskip
\bigskip
\textbf{Question 5.7}
\bigskip

An investment will pay \$150 at the end of each of the next 3 years, \$250 at the end of year 4, \$300 at the end of year 5, and \$600 at the end of year 6. 
If other investments of equal risk earn 11\% annually, what is its present value and its future value?

\bigskip
\textbf{Answer}

\[
  PV = \frac {150}{1.11} + \frac {150}{1.11^2} + \frac {150}{1.11^3} + \frac {250}{1.11^4} + \frac {300}{1.11^5} + \frac {600}{1.11^6}
\]

\[
  \boxed{PV \approx \$1926.64}
\]

\bigskip
\bigskip
\textbf{Question 5.23}
\bigskip

Find the amount to which \$500 will grow under each of these conditions:

\begin{enumerate}
  \item[a.] 12\% compounded annually for 5 years
  \item[b.] 12\% compounded semiannually for 5 years
  \item[c.] 12\% compounded quarterly for 5 years
  \item[d.] 12\% compounded monthly for 5 years
  \item[e.] 12\% compounded daily for 5 years
  \item[f.] Why does the observed pattern of FVs occur? 
\end{enumerate}

\bigskip
\textbf{Answer}

\begin{enumerate}
  \item[a.] $500(1.12)^5 \approx 881.17$
  \item[b.] $5000(1 + \frac {0.12}{2})^{5(2)} \approx 895.42$ 
  \item[c.] $5000(1 + \frac {0.12}{4})^{5(4)} \approx 903.06$ 
  \item[d.] $5000(1 + \frac {0.12}{12})^{5(12)} \approx 908.35$ 
  \item[e.] $5000(1 + \frac {0.12}{365})^{5(365)} \approx 911.00$ 
  \item[f.] Shorter payment periods result in more frequent compounding, thus increasing total value. 
\end{enumerate}

\end{document}